\section*{Hoofdstuk \ref{ch:g1:the-five-senses} --- De vijf zintuigen}

\begin{solution}{exo:g1:the-five-senses:1}
Zien --- de \omterm{def:g1:the-five-senses:sense}{ogen}; horen --- de \omterm{def:g1:the-five-senses:sense}{oren}; ruiken --- de \omterm{def:g1:the-five-senses:sense}{neus}; proeven ---
de \omterm{def:g1:the-five-senses:sense}{tong}; voelen --- de \omterm{def:g1:the-five-senses:sense}{huid}.
\end{solution}

\begin{solution}{exo:g1:the-five-senses:2}
(a) De \omterm{def:g1:the-five-senses:sense}{ogen}. (b) De \omterm{def:g1:the-five-senses:sense}{huid}. (c) De \omterm{def:g1:the-five-senses:sense}{neus}. (d) De \omterm{def:g1:the-five-senses:sense}{oren}.
\end{solution}

\begin{solution}{exo:g1:the-five-senses:3}
De \omterm{def:g1:the-five-senses:sense}{huid}, het orgaan van het voelen. Ze bedekt het hele lichaam, van top
tot teen.
\end{solution}

\begin{solution}{exo:g1:the-five-senses:4}
De \omterm{def:g1:the-five-senses:sense}{huid} van de handen stuurde de boodschap (koud, hard, gekarteld); de
hersenen begrepen hem en noemden het voorwerp.
\end{solution}

\begin{solution}{exo:g1:the-five-senses:5}
Proeven en ruiken. De \omterm{def:g1:the-five-senses:sense}{neus} doet de helft van het werk van de smaak, en
daarom lijkt eten flauw als hij verstopt is.
\end{solution}

\begin{solution}{exo:g1:the-five-senses:6}
Bijvoorbeeld: kijk nooit recht in de Zon --- beschermt de \omterm{def:g1:the-five-senses:sense}{ogen}; houd
geluiden zacht --- beschermt de \omterm{def:g1:the-five-senses:sense}{oren}; blaas op heet eten --- beschermt
de \omterm{def:g1:the-five-senses:sense}{tong}; niets in \omterm{def:g1:the-five-senses:sense}{oor} of \omterm{def:g1:the-five-senses:sense}{neus} duwen --- beschermt \omterm{def:g1:the-five-senses:sense}{oren} en \omterm{def:g1:the-five-senses:sense}{neus}; goed
licht om te lezen --- beschermt de \omterm{def:g1:the-five-senses:sense}{ogen}.
\end{solution}

\begin{solution}{exo:g1:the-five-senses:7}
Omdat te hard geluid de \omterm{def:g1:the-five-senses:sense}{oren} beschadigt, en een beschadigd \omterm{def:g1:the-five-senses:sense}{oor}
beschadigd kan blijven. De beschermers houden het lawaai van de
machines zacht genoeg voor een leven lang horen.
\end{solution}

\begin{solution}{exo:g1:the-five-senses:8}
Open antwoord --- de meeste mensen noemen er twee of drie van de drie.
De \omterm{def:g1:the-five-senses:sense}{neus} herkent bekend eten goed, zelfs zonder hulp van de \omterm{def:g1:the-five-senses:sense}{ogen}.
\end{solution}

\begin{solution}{exo:g1:the-five-senses:9}
Het \omterm{def:g1:the-five-senses:sense}{oog} \emph{ziet} alleen --- het vangt licht op en stuurt een
boodschap. De hersenen doen het begrijpen: ze ontvangen de boodschap en
zeggen ``dat is oma'', ``dat is mijn school''. \omterm{def:g1:the-five-senses:sense}{Zintuigen} verzamelen; de
hersenen begrijpen.
\end{solution}

\begin{solution}{exo:g1:the-five-senses:10}
Lezen: met de vingertoppen --- het gevoel neemt het van het zien over
en voelt verhoogde puntjes op de bladzijde. Oversteken: met de \omterm{def:g1:the-five-senses:sense}{oren} ---
luisteren naar auto's die komen en gaan --- vaak met een stok waarvan
de punt de grond vooruit laat aftasten.
\end{solution}
